\begingroup\let\clearpageforsection\relax\exercisesheader{}\endgroup

% 13

\eoce{\qt{Pencemaran udara dan luaran kelahiran: cakupan inferensi\label{scope_airpoll}} 
Latihan~\ref{study_components_airpoll} memperkenalkan sebuah studi yang 
mengumpulkan data untuk menelaah hubungan antara pencemar udara dan kelahiran 
prematur di California Selatan. Selama studi, tingkat pencemaran udara diukur 
oleh stasiun pemantau kualitas udara. Data lama kehamilan dikumpulkan untuk 
143.196 kelahiran antara 1989 dan 1993, dan paparan pencemaran udara selama 
kehamilan dihitung untuk setiap kelahiran.
\begin{parts}
\item Tentukan populasi sasaran dan sampel dalam studi ini.
\item Jelaskan apakah hasil studi dapat digeneralisasikan ke populasi dan apakah 
temuannya dapat digunakan untuk menetapkan hubungan kausal.
\end{parts}
}{}

% 14

\eoce{\qt{Kecurangan: cakupan inferensi\label{scope_cheaters}} 
Latihan~\ref{study_components_cheaters} memperkenalkan studi tentang hubungan 
antara kejujuran, usia, dan pengendalian diri. Para peneliti melakukan eksperimen 
pada 160 anak berusia 5 hingga 15 tahun. Setiap anak diminta melempar koin 
seimbang tanpa dilihat orang lain dan mencatat hasilnya (putih atau hitam) pada 
selembar kertas; hanya anak yang melaporkan hasil putih yang akan diberi hadiah. 
Separuh anak 
secara eksplisit diminta untuk tidak berbuat curang, sedangkan separuh lainnya 
tidak diberi instruksi eksplisit. Terdapat perbedaan tingkat kecurangan antara 
kelompok dengan dan tanpa instruksi, serta sejumlah perbedaan menurut 
karakteristik anak di dalam setiap kelompok.
\begin{parts}
\item Tentukan populasi sasaran dan sampel dalam studi ini.
\item Jelaskan apakah hasil studi dapat digeneralisasikan ke populasi dan apakah 
temuannya dapat digunakan untuk menetapkan hubungan kausal.
\end{parts}
}{}

% 15

\eoce{\qt{Metode Buteyko: cakupan inferensi\label{scope_buteyko}} 
Latihan~\ref{study_components_buteyko} memperkenalkan studi mengenai penggunaan 
teknik pernapasan dangkal Buteyko untuk mengurangi gejala asma dan meningkatkan 
kualitas hidup. Studi ini merekrut 600 pasien asma berusia 18--69 tahun yang 
mengandalkan obat untuk perawatan asma, lalu mengalokasikan mereka secara acak 
ke dua kelompok: satu kelompok mempraktikkan metode Buteyko dan kelompok lainnya 
tidak. Kelompok Buteyko rata-rata mengalami penurunan gejala asma yang signifikan 
dan peningkatan kualitas hidup.
\begin{parts}
\item Tentukan populasi sasaran dan sampel dalam studi ini.
\item Jelaskan apakah hasil studi dapat digeneralisasikan ke populasi dan apakah 
temuannya dapat digunakan untuk menetapkan hubungan kausal.
\end{parts}
}{}

% 16

\eoce{\qt{Mengambil permen: cakupan inferensi\label{scope_stealers}} 
Latihan~\ref{study_components_stealers} memperkenalkan studi tentang hubungan 
antara kelas sosial ekonomi dan perilaku tidak etis. Dalam studi ini, 129 
mahasiswa program sarjana University of California, Berkeley diminta 
menggolongkan diri ke kelas sosial rendah atau tinggi dengan membandingkan diri 
mereka dengan orang yang memiliki paling banyak (atau paling sedikit) uang dan 
pendidikan, serta pekerjaan yang paling (atau paling tidak) dihormati. Mereka 
juga diberi sebuah stoples berisi permen yang dibungkus satu per satu. Peneliti 
memberi tahu bahwa permen tersebut ditujukan bagi anak-anak di laboratorium 
terdekat, tetapi peserta boleh mengambil beberapa jika ingin. Setelah 
menyelesaikan tugas lain yang tidak berkaitan, peserta melaporkan jumlah permen 
yang mereka ambil. Peserta yang tergolong kelas atas ditemukan mengambil lebih 
banyak permen daripada peserta lainnya.
\begin{parts}
\item Tentukan populasi sasaran dan sampel dalam studi ini.
\item Jelaskan apakah hasil studi dapat digeneralisasikan ke populasi dan apakah 
temuannya dapat digunakan untuk menetapkan hubungan kausal.
\end{parts}
}{}

% 17

\eoce{\qt{Bersantai setelah bekerja\label{relax_after_work_definitions}} General 
Social Survey mengambil sampel acak 1.155 warga Amerika dan menanyakan berapa 
jam yang mereka miliki setelah hari kerja biasa untuk bersantai atau melakukan 
kegiatan yang mereka sukai. Rata-rata waktu bersantai yang diperoleh adalah 
1,65 jam. Tentukan mana di antara hal-hal berikut yang merupakan observasi, 
variabel, statistik sampel (nilai yang dihitung berdasarkan sampel teramati), 
atau parameter populasi.
\begin{parts}
\item Seorang warga Amerika dalam sampel.
\item Jumlah jam yang digunakan untuk bersantai setelah hari kerja biasa.
\item 1,65.
\item Rata-rata jumlah jam yang digunakan seluruh warga Amerika untuk bersantai 
setelah hari kerja biasa.
\end{parts}
}{}

% 18

\eoce{\qt{Kucing di YouTube\label{cats_on_youtube_definitions}} Misalkan Anda 
ingin memperkirakan persentase video YouTube yang menampilkan kucing. Karena 
mustahil menonton semua video di YouTube, Anda menggunakan pemilih video acak 
untuk memilih 1.000 video. Anda menemukan bahwa 2\% dari video tersebut 
menampilkan kucing. Tentukan mana di antara hal-hal berikut yang merupakan 
observasi, variabel, statistik sampel (nilai yang dihitung berdasarkan sampel 
teramati), atau parameter populasi.
\begin{parts}
\item Persentase seluruh video YouTube yang menampilkan kucing.
\item 2\%.
\item Sebuah video dalam sampel Anda.
\item Apakah sebuah video menampilkan kucing atau tidak.
\end{parts}
}{}

% 19

\eoce{\qt{Kepuasan mata kuliah antarkelompok\label{course_satisfaction_sections}} 
Sebuah kelas besar di perguruan tinggi memiliki 160 mahasiswa. Seluruh 
mahasiswa mengikuti kuliah bersama, tetapi untuk sesi praktikum mereka dibagi 
menjadi 4 kelompok yang masing-masing berisi 40 mahasiswa dan dibimbing oleh 
asisten pengajar yang berbeda. Dosen hendak melakukan survei mengenai kepuasan 
mahasiswa terhadap mata kuliah tersebut. Ia menduga kelompok praktikum yang 
diikuti mahasiswa dapat memengaruhi kepuasan mereka secara keseluruhan.
\begin{parts}
\item Jenis studi apakah ini?
\item Usulkan strategi pengambilan sampel untuk melaksanakan studi ini.
\end{parts}
}{}

% 20

\eoce{\qt{Usulan hunian di berbagai asrama\label{housing_proposal_dorms}} Di sebuah 
kampus besar, mahasiswa tahun pertama dan kedua tinggal di asrama di bagian 
timur kampus, sedangkan mahasiswa tahun ketiga dan keempat tinggal di asrama 
di bagian barat. Misalkan Anda ingin mengumpulkan pendapat mahasiswa tentang 
struktur hunian baru yang diusulkan pengelola perguruan tinggi dan ingin 
memastikan bahwa survei mewakili pendapat mahasiswa dari setiap angkatan 
secara setara.
\begin{parts}
\item Jenis studi apakah ini?
\item Usulkan strategi pengambilan sampel untuk melaksanakan studi ini.
\end{parts}
}{}

% 21

\eoce{\qt{Penggunaan internet dan harapan hidup\label{internet_life_expectancy}} 
Diagram pencar berikut dibuat dalam sebuah studi yang menilai hubungan antara 
estimasi harapan hidup saat lahir (per 2014) dan persentase pengguna internet 
(per 2009) di 208 negara yang memiliki kedua jenis data tersebut.\footfullcite{data:ciaFactbook}

\noindent\begin{minipage}[c]{0.44\textwidth}
\begin{parts}
\item Jelaskan hubungan antara harapan hidup dan persentase pengguna internet.
\item Jenis studi apakah ini?
\item Sebutkan kemungkinan variabel perancu yang dapat menjelaskan hubungan ini 
dan uraikan kemungkinan pengaruhnya.
\end{parts} \vspace{15mm}
\end{minipage}
\begin{minipage}[r]{0.55\textwidth}
\hfill%
\Figures[Diagram pencar dengan "persentase pengguna internet" (0\% hingga 100\%) pada sumbu horizontal dan "harapan hidup saat lahir" (50 hingga 90) pada sumbu vertikal. Untuk 0\% hingga 15\%, sekitar 100 titik tersebar cukup merata antara 50 dan 75. Untuk 15\% hingga 90\%, titik-titik terkonsentrasi antara sekitar 70 dan 85 serta menunjukkan sedikit tren menaik.]{0.87}{eoce/internet_life_expectancy}{internet_life_expectancy}
\end{minipage}
}{}

% 22

\eoce{\qt{Dilanda stres, Bagian I\label{stressed_out_observational}} Sebuah studi 
yang menyurvei sampel acak siswa sekolah menengah yang secara umum sehat 
menemukan bahwa mereka lebih mungkin mengalami kram otot ketika stres. Studi 
tersebut juga mencatat bahwa siswa minum lebih banyak kopi dan tidur lebih 
sedikit ketika stres.
\begin{parts}
\item Jenis studi apakah ini?
\item Dapatkah studi ini digunakan untuk menyimpulkan hubungan kausal antara 
peningkatan stres dan kram otot?
\item Sebutkan kemungkinan variabel perancu yang dapat menjelaskan hubungan 
teramati antara peningkatan stres dan kram otot. 
\end{parts}
}{}

% 23

\eoce{\qt{Evaluasi metode pengambilan sampel\label{evaluate_sampling_methods}} 
Sebuah universitas ingin mengetahui proporsi mahasiswa program sarjananya yang 
mendukung iuran tahunan baru sebesar \$25 untuk memperbaiki pusat kegiatan 
mahasiswa. Untuk setiap metode yang diusulkan berikut, nyatakan apakah metode 
tersebut masuk akal atau tidak.
\begin{parts}
\item Survei sampel acak sederhana yang terdiri atas 500 mahasiswa.
\item Bagi mahasiswa ke dalam strata berdasarkan bidang studi, lalu ambil sampel 
10\% mahasiswa dari setiap strata.
\item Kelompokkan mahasiswa ke dalam klaster berdasarkan usia (misalnya usia 
18 tahun dalam satu klaster, usia 19 tahun dalam satu klaster, dan seterusnya), 
lalu pilih tiga klaster secara acak dan survei seluruh mahasiswa dalam klaster 
tersebut.
\end{parts}
}{}

% 24

\eoce{\qt{Pemanggilan digit acak\label{random_digit_dialing}} Gallup Poll 
menggunakan prosedur yang disebut pemanggilan digit acak. Prosedur ini membuat 
nomor telepon berdasarkan daftar seluruh kode area di Amerika Serikat serta 
jumlah rumah tangga yang terkait dengan setiap kode area. Berikan satu alasan 
yang mungkin menjelaskan mengapa Gallup Poll menggunakan pemanggilan digit acak 
alih-alih memilih nomor telepon dari buku telepon.
}{}

% 25

\eoce{\qt{Kecenderungan menyukai atau tidak menyukai\label{scope_haters}}
Sebuah studi yang diterbitkan dalam
\textit{Journal of Personality and Social Psychology}
meminta 200 peserta dari kalangan laki-laki dan perempuan yang dipilih secara
acak untuk menilai
perasaan mereka terhadap berbagai topik, seperti berkemah, layanan kesehatan,
arsitektur, taksidermi, teka-teki silang, dan Jepang. Tujuannya adalah mengukur
sikap mereka terhadap rangsangan yang sebagian besar tidak saling berkaitan.
Selanjutnya, peserta menerima informasi tentang produk baru berupa oven
gelombang mikro. Produk ini sebenarnya tidak ada, tetapi peserta tidak
mengetahuinya; mereka diberi tiga ulasan positif dan tiga ulasan negatif yang
direkayasa. Orang yang bereaksi positif terhadap topik-topik pada pengukuran
sikap disposisional juga cenderung bereaksi positif terhadap oven tersebut,
sedangkan orang yang bereaksi negatif cenderung menilainya secara negatif.
Para peneliti menafsirkan pola itu sebagai kecenderungan evaluatif umum:
sebagian peserta cenderung menilai banyak objek secara lebih positif, sedangkan
peserta lain cenderung menilainya secara lebih negatif. Menurut mereka, pola
lintas-objek ini dapat membantu menjelaskan bagaimana sikap terbentuk.
\footfullcite{Hepler:2013}
\begin{parts}
\item Apa saja yang menjadi kasus?
\item Apa variabel respons dalam studi ini?
\item Apa variabel penjelas dalam studi ini?
\item Apakah studi ini menggunakan pengambilan sampel acak?
\item Apakah ini studi observasional atau eksperimen? Jelaskan alasan Anda.
\item Dapatkah kita menetapkan hubungan kausal antara variabel penjelas dan 
variabel respons?
\item Dapatkah hasil studi digeneralisasikan ke populasi secara luas?
\end{parts}
}{}

% 26

\eoce{\qt{Ukuran keluarga\label{family_size}} Misalkan kita ingin memperkirakan 
ukuran rumah tangga. Dalam konteks ini, rumah tangga terdiri atas orang-orang 
yang tinggal bersama dalam hunian yang sama dan berbagi fasilitas tempat 
tinggal. Jika kita memilih murid secara acak di sebuah sekolah dasar dan 
menanyakan ukuran keluarga mereka, apakah hasilnya akan mengukur ukuran rumah 
tangga dengan baik? Atau apakah rata-ratanya akan bias? Jika bias, apakah 
hasilnya akan mengestimasi nilai sebenarnya terlalu tinggi atau terlalu rendah?
}{}

% 27

\eoce{\qt{Strategi pengambilan sampel\label{sampling_strategies}} Seorang mahasiswa 
statistika ingin mengetahui hubungan antara waktu yang dihabiskan siswa di situs 
jejaring sosial dan kinerja mereka di sekolah, lalu memutuskan untuk melakukan 
survei. Beberapa strategi penelitian untuk mengumpulkan data diuraikan berikut. 
Untuk setiap strategi, sebutkan metode pengambilan sampel yang diusulkan dan 
bias yang mungkin timbul.
\begin{parts}
\item Ia memilih 40 siswa secara acak dari populasi studi, memberikan survei 
kepada mereka, lalu meminta mereka mengisinya dan mengembalikannya keesokan hari.
\item Ia hanya membagikan survei kepada teman-temannya dan memastikan setiap 
teman mengisinya.
\item Ia mengunggah tautan survei daring di Facebook dan meminta teman-temannya 
mengisi survei tersebut.
\item Ia memilih 5 kelas secara acak, lalu meminta sampel acak siswa dari 
kelas-kelas tersebut untuk mengisi survei.
\end{parts}
}{}

% 28

\eoce{\qt{Membaca berita\label{reading_paper}} Dua artikel di \emph{NY Times}
memberitakan studi berikut. Uraian di bawah adalah ringkasan faktual yang
disusun ulang, bukan kutipan dari artikel:
\begin{parts}
\item Artikel berjudul \emph{Risks: Smokers Found More Prone to Dementia} 
melaporkan hal-hal berikut. \footfullcite{news:smokingDementia}
\begin{adjustwidth}{1em}{1em}
{\footnotesize Analisis mencakup 23.123 anggota program kesehatan yang berusia
50--60 tahun ketika secara sukarela menjalani pemeriksaan dan survei perilaku
pada 1978--1985. Setelah 23 tahun, sekitar 25\% peserta tercatat mengalami
demensia; 1.136 kasus adalah penyakit Alzheimer dan 416 kasus adalah demensia
vaskular. Dalam perbandingan yang telah disesuaikan terhadap faktor lain,
perkiraan risiko relatif terhadap bukan perokok adalah 37\% lebih tinggi untuk
satu bungkus per hari, 44\% lebih tinggi untuk satu hingga dua bungkus, dan
sekitar dua kali lipat untuk lebih dari dua bungkus per hari.}
\end{adjustwidth}
Berdasarkan studi ini, dapatkah kita menyimpulkan bahwa merokok menyebabkan 
demensia di kemudian hari? Jelaskan alasan Anda.
\item Artikel lain berjudul \emph{The School Bully Is Sleepy} melaporkan hal-hal 
berikut. \footfullcite{news:bullySleep}
\begin{adjustwidth}{1em}{1em}
{\footnotesize Studi University of Michigan menggabungkan laporan orang tua
tentang kebiasaan tidur anak dengan penilaian orang tua dan guru mengenai
perilaku. Sekitar sepertiga siswa memperoleh penandaan masalah perilaku
mengganggu atau perundungan. Gejala gangguan tidur dilaporkan sekitar dua kali
lebih sering pada anak dalam kelompok bermasalah tersebut dan pada anak yang
diidentifikasi sebagai perundung.}
\end{adjustwidth}
Seorang teman yang membaca artikel itu menyimpulkan bahwa gangguan tidur 
menyebabkan perundungan pada anak sekolah. Apakah pernyataan tersebut dapat 
dibenarkan? Jika tidak, bagaimana sebaiknya Anda menjelaskan kesimpulan yang 
dapat ditarik dari studi ini?
\end{parts}
}{}
